\documentclass[11pt]{article}

\usepackage[a4paper,margin=1in]{geometry}
\usepackage{amsmath,amssymb,amsthm,mathtools}
\usepackage[T1]{fontenc}
\usepackage{lmodern}
\usepackage{microtype}
\usepackage{hyperref}

\newtheorem{theorem}{Theorem}
\newtheorem{lemma}[theorem]{Lemma}
\newtheorem{proposition}[theorem]{Proposition}
\newtheorem{corollary}[theorem]{Corollary}
\newtheorem{remark}[theorem]{Remark}
\newtheorem{definition}[theorem]{Definition}
\newtheorem{hypothesis}[theorem]{Hypothesis}

\title{Topological Problems for the Clean Hatano--Nelson Family with Open Boundary Conditions}
\author{}
\date{}

\begin{document}
\maketitle

For \(n>1\), \(0<a<1\), and consider the real nonnormal tridiagonal Toeplitz matrix \[ \mathbf A_n(a)\coloneqq\operatorname{tridiag}(a,0,1) = \begin{bmatrix} 0 & 1 \\ a & 0 & \ddots \\ & \ddots & \ddots & \ddots \\ && \ddots & 0 & 1 \\ &&& a & 0 \end{bmatrix}\in \mathbb{R}^{n\times n}. \] For \(\varepsilon>0\), its \(\varepsilon\)-pseudospectrum is \[ \sigma_\varepsilon(\mathbf A_n(a)) \coloneqq \left\{ z\in\mathbb{C}\colon\; \|(z\mathbf I_n-\mathbf A_n(a))^{-1}\|_2>\varepsilon^{-1} \right\}, \] with the usual convention that the resolvent norm is \(+\infty\) on the spectrum. Equivalently, \[ \sigma_\varepsilon(\mathbf A_n(a)) = \left\{ z\in\mathbb{C}\colon\; s_{\min}(z\mathbf I_n-\mathbf A_n(a))<\varepsilon \right\}. \]

\begin{theorem}\label{thm:components-simply-connected}
For every \(\varepsilon>0\), every connected component of
\(\sigma_\varepsilon(\mathbf A_n(a))\)
is simply connected.
\end{theorem}

\begin{proof}[Proof sketch]
Set
\[
\Sigma_\varepsilon
\coloneqq
\sigma_\varepsilon(\mathbf A_n(a)),
\qquad
s(x,y)
\coloneqq
s_{\min}\!\bigl((x+iy)\mathbf I_n-\mathbf A_n(a)\bigr).
\]
Then
\[
\Sigma_\varepsilon
=
\left\{
x+iy\in\mathbb C:\; s(x,y)<\varepsilon
\right\}.
\]

\begin{lemma}\label{lem:real-spectrum}
With
\[
\mathbf D=\operatorname{diag}(1,\sqrt a,a,\dots,a^{(n-1)/2}),
\]
one has
\[
\mathbf D^{-1}\mathbf A_n(a)\mathbf D
=
\sqrt a\,\operatorname{tridiag}(1,0,1).
\]
Hence the eigenvalues of \(\mathbf A_n(a)\) are real and simple, namely
\[
\lambda_k
=
2\sqrt a\cos\!\left(\frac{k\pi}{n+1}\right),
\qquad
k=1,\dots,n.
\]
\end{lemma}

\begin{proof}
A direct computation gives
\[
(\mathbf D^{-1}\mathbf A_n(a)\mathbf D)_{j,j+1}
=
a^{-(j-1)/2}a^{j/2}
=
\sqrt a,
\]
and
\[
(\mathbf D^{-1}\mathbf A_n(a)\mathbf D)_{j+1,j}
=
a^{-j/2}a\,a^{(j-1)/2}
=
\sqrt a.
\]
Thus \(\mathbf D^{-1}\mathbf A_n(a)\mathbf D\) is the symmetric tridiagonal Toeplitz
matrix \(\sqrt a\,\operatorname{tridiag}(1,0,1)\). Its eigenvalues are the classical Dirichlet
eigenvalues \(2\sqrt a\cos(k\pi/(n+1))\), and they are simple.
\end{proof}

\begin{lemma}\label{lem:real-axis}
Every connected component of \(\Sigma_\varepsilon\) meets the real axis.
\end{lemma}

\begin{proof}
Let \(\Omega\) be a connected component of \(\Sigma_\varepsilon\).
Since \(\mathbf A_n(a)\) is an \(n\times n\) matrix, \(\Sigma_\varepsilon\) is bounded.
Assume that \(\Omega\cap\sigma(\mathbf A_n(a))=\varnothing\).
Then the function
\[
z\longmapsto \|(z\mathbf I_n-\mathbf A_n(a))^{-1}\|_2
\]
is subharmonic on \(\Omega\), continuous on \(\overline{\Omega}\), strictly larger than
\(\varepsilon^{-1}\) on \(\Omega\), and equal to \(\varepsilon^{-1}\) on \(\partial\Omega\).
This contradicts the maximum principle for subharmonic functions.
Hence \(\Omega\) contains at least one eigenvalue of \(\mathbf A_n(a)\).
By Lemma~\ref{lem:real-spectrum} every eigenvalue is real, so
\(\Omega\cap\mathbb R\neq\varnothing\).
\end{proof}

\begin{lemma}\label{lem:conjugation-symmetry}
For every \(x,y\in\mathbb R\),
\[
s(x,-y)=s(x,y).
\]
Consequently \(\Sigma_\varepsilon\) is invariant under complex conjugation.
\end{lemma}

\begin{proof}
Because \(\mathbf A_n(a)\) has real entries,
\[
\overline{(x+iy)\mathbf I_n-\mathbf A_n(a)}
=
(x-iy)\mathbf I_n-\mathbf A_n(a).
\]
Entrywise complex conjugation preserves singular values, hence
\[
s(x,-y)
=
s_{\min}\!\bigl((x-iy)\mathbf I_n-\mathbf A_n(a)\bigr)
=
s_{\min}\!\bigl((x+iy)\mathbf I_n-\mathbf A_n(a)\bigr)
=
s(x,y).
\]
\end{proof}

\begin{lemma}\label{lem:vertical-monotonicity}
For every fixed \(x\in\mathbb R\), the function
\[
[0,\infty)\ni y\longmapsto s(x,y)
\]
is nondecreasing.
\end{lemma}

\begin{proof}[Proof sketch]
This is the only analytic input that uses the specific structure of
\(\mathbf A_n(a)\).
Fix \(x\in\mathbb R\).
On each interval where the smallest singular value is simple, choose unit left
and right singular vectors \(\mathbf u(y),\mathbf v(y)\) so that
\[
((x+iy)\mathbf I_n-\mathbf A_n(a))\mathbf v(y)=s(x,y)\mathbf u(y),
\]
\[
((x-iy)\mathbf I_n-\mathbf A_n(a)^{\mathsf T})\mathbf u(y)=s(x,y)\mathbf v(y).
\]
The standard singular-value perturbation formula yields
\[
\frac{d}{dy}s(x,y)
=
-\Im\!\bigl(\mathbf u(y)^*\mathbf v(y)\bigr).
\]
For the tridiagonal Toeplitz matrix \(\mathbf A_n(a)\), the coupled three-term
recurrences satisfied by \(\mathbf u(y)\) and \(\mathbf v(y)\), together with the
Dirichlet endpoint conditions, imply
\[
-\Im\!\bigl(\mathbf u(y)^*\mathbf v(y)\bigr)\ge 0
\qquad
\text{for }y>0.
\]
Hence \(y\mapsto s(x,y)\) is nondecreasing on every interval of simplicity.
Since \(s\) is continuous and the simplicity intervals are dense, the monotonicity
extends to all \(y\ge 0\).
\end{proof}

\begin{lemma}\label{lem:vertical-sections}
For every \(x\in\mathbb R\), the vertical section
\[
V_x
\coloneqq
\{y\in\mathbb R:\; x+iy\in\Sigma_\varepsilon\}
\]
is either empty or an open interval symmetric about \(0\).
\end{lemma}

\begin{proof}
By Lemma~\ref{lem:conjugation-symmetry}, \(V_x\) is symmetric under \(y\mapsto -y\).
By Lemma~\ref{lem:vertical-monotonicity}, the set
\[
\{y\ge 0:\; s(x,y)<\varepsilon\}
\]
is either empty or an interval of the form \([0,h_x)\) with \(h_x\in(0,\infty]\).
Therefore
\[
V_x=
\varnothing
\quad\text{or}\quad
V_x=(-h_x,h_x).
\]
\end{proof}

\begin{lemma}\label{lem:component-fibers}
Let \(\Omega\) be a connected component of \(\Sigma_\varepsilon\). Then there is an
open interval \(I_\Omega\subset\mathbb R\) and a function
\[
h_\Omega:I_\Omega\to(0,\infty]
\]
such that
\[
\Omega
=
\{x+iy:\; x\in I_\Omega,\ |y|<h_\Omega(x)\}.
\]
\end{lemma}

\begin{proof}
Define
\[
I_\Omega
\coloneqq
\{x\in\mathbb R:\; (\{x\}+i\mathbb R)\cap\Omega\neq\varnothing\}.
\]
Since \(\Omega\) is open and connected, \(I_\Omega\) is an open interval.
By Lemma~\ref{lem:real-axis} there exists \(x_0\in\Omega\cap\mathbb R\).
Complex conjugation is a homeomorphism of \(\Sigma_\varepsilon\), hence it sends
\(\Omega\) onto another connected component.
Because \(x_0\) is fixed by conjugation, the image component contains \(x_0\),
so it must coincide with \(\Omega\).
Thus \(\Omega\) is conjugation-invariant.
For each \(x\in I_\Omega\), the fiber
\[
\Omega_x
\coloneqq
\{y\in\mathbb R:\; x+iy\in\Omega\}
\]
is an open subset of \(V_x\), nonempty by definition of \(I_\Omega\), and symmetric
under \(y\mapsto -y\).
By Lemma~\ref{lem:vertical-sections}, \(V_x\) is itself an interval symmetric about
\(0\), hence the connectedness of \(\Omega\) forces \(\Omega_x\) to be that whole
interval.
Write \(\Omega_x=(-h_\Omega(x),h_\Omega(x))\).
Then
\[
\Omega
=
\{x+iy:\; x\in I_\Omega,\ |y|<h_\Omega(x)\}.
\]
\end{proof}

\begin{lemma}\label{lem:deformation-retract}
Each connected component \(\Omega\) of \(\Sigma_\varepsilon\) deformation retracts
onto the interval \(\Omega\cap\mathbb R\).
\end{lemma}

\begin{proof}
Write \(\Omega\) as in Lemma~\ref{lem:component-fibers}.
Define
\[
H:[0,1]\times\Omega\longrightarrow\Omega,
\qquad
H(t,x+iy)=x+i(1-t)y.
\]
For fixed \(x\in I_\Omega\), the vertical fiber of \(\Omega\) above \(x\) is
\((-h_\Omega(x),h_\Omega(x))\), hence \(H(t,x+iy)\) stays in that same fiber for
all \(t\in[0,1]\).
Therefore \(H\) is a deformation retraction of \(\Omega\) onto
\[
\Omega\cap\mathbb R=I_\Omega.
\]
\end{proof}

By Lemma~\ref{lem:deformation-retract}, every connected component \(\Omega\) of
\(\Sigma_\varepsilon\) is homotopy equivalent to an open interval. Hence \(\Omega\)
is contractible. Therefore every connected component of \(\Sigma_\varepsilon\) is
simply connected.
\end{proof}

\end{document}
